\chapter{\abstractname}

Wearables and small drones need to understand how they move while operating
with limited power and computation. Visual odometry estimates this movement
from camera images, usually by tracking image features over time. Yet many
camera devices already contain dedicated hardware that searches for matching
image regions to compress video. Could the motion information produced by
this widely available hardware also support camera tracking?

This thesis explores that possibility using motion vectors from H.264 video
compression. These vectors describe how the encoder predicts image blocks;
they are designed to save data, so they do not necessarily follow the same
physical points over time. The proposed method combines vectors from two
cameras with known camera geometry to estimate depth and camera movement.
Its tracking and motion estimation use these measurements without consulting
image pixels. This approach is compatible with compress-then-analyze: a
camera could send the required compressed streams to a companion phone,
which reuses their motion information for localisation.

Experiments on an established drone benchmark recover coherent camera
trajectories, particularly on easier sequences. More demanding sequences
expose limitations, and methods that refine their tracks using image pixels
remain more accurate overall. The results nevertheless demonstrate that
useful camera tracking from codec motion vectors is feasible. They provide
one starting point for turning existing video-processing resources into a
source of geometric information. The experiments establish this possibility;
live hardware operation and the potential energy savings remain to be
evaluated.


\makeatletter
\ifthenelse{\pdf@strcmp{\languagename}{english}=0}
{\renewcommand{\abstractname}{Kurzfassung}}
{\renewcommand{\abstractname}{Abstract}}
\makeatother

\chapter{\abstractname}

\begin{otherlanguage}{ngerman}
\sisetup{locale=DE}

Wearables und kleine Drohnen müssen ihre eigene Bewegung erfassen können,
obwohl ihnen nur begrenzte Energie und Rechenleistung zur Verfügung stehen.
Visuelle Odometrie schätzt diese Bewegung aus Kamerabildern, üblicherweise
durch die Verfolgung von Bildmerkmalen über die Zeit. Viele Kamerageräte
enthalten jedoch bereits spezialisierte Hardware, die zur Videokompression
nach übereinstimmenden Bildregionen sucht. Könnten die dabei gewonnenen
Bewegungsinformationen auch zur Schätzung der Kamerabewegung dienen?

Diese Arbeit untersucht diese Möglichkeit anhand von Bewegungsvektoren
aus der H.264-Videokompression. Die Vektoren beschreiben, wie der Encoder
Bildblöcke vorhersagt. Da sie der Datenreduktion dienen, verfolgen sie nicht
zwangsläufig dieselben physischen Punkte über die Zeit. Das vorgeschlagene
Verfahren kombiniert Vektoren aus zwei Kameras mit bekannter Kamerageometrie,
um Tiefe und Kamerabewegung zu schätzen. Die Verfolgung und Bewegungsschätzung
nutzen diese Messungen, ohne auf Bildpixel zuzugreifen. Dieser Ansatz ist
mit Compress-then-Analyze vereinbar: Eine Kamera könnte die benötigten
komprimierten Videoströme an ein gekoppeltes Smartphone senden, das deren
Bewegungsinformationen für die Lokalisierung wiederverwendet.

Experimente auf einem etablierten Drohnen-Benchmark liefern schlüssige
Kameratrajektorien, insbesondere auf einfacheren Sequenzen. Anspruchsvollere
Sequenzen zeigen Grenzen auf; Verfahren, die ihre Korrespondenzen anhand von
Bildpixeln verfeinern, bleiben insgesamt genauer. Dennoch belegen die
Ergebnisse, dass eine brauchbare Schätzung der Kamerabewegung allein aus
Codec-Bewegungsvektoren möglich ist. Sie bieten einen Ausgangspunkt dafür,
vorhandene Ressourcen zur Videoverarbeitung als Quelle geometrischer
Informationen zu nutzen. Die Experimente zeigen diese Machbarkeit; der
Live-Betrieb auf entsprechender Hardware und mögliche Energieeinsparungen
müssen noch untersucht werden.

\end{otherlanguage}


% Undo the name switch
\makeatletter
\ifthenelse{\pdf@strcmp{\languagename}{english}=0}
{\renewcommand{\abstractname}{Abstract}}
{\renewcommand{\abstractname}{Kurzfassung}}
\makeatother
